\documentclass{beamer}

\usepackage{beamerthemeSingapore}
\usepackage[brazil]{babel}
\usepackage[T1]{fontenc}
\usepackage[utf8]{inputenc}
\usepackage{graphicx}
\usepackage{url}
\usepackage{xcolor}
\usepackage{booktabs}
\usepackage{ragged2e}
\usepackage{etoolbox}
\usepackage{minted}

\newcommand{\tcb}[1]{\textcolor{blue}{#1}}
\newcommand{\tcr}[1]{\textcolor{red}{#1}}
\newcommand{\tcg}[1]{\textcolor{green!50!black}{#1}}

\setminted{
  fontsize=\footnotesize,
  breaklines=true,
  breakanywhere=true,
  frame=lines,
  framesep=2mm,
  baselinestretch=1.05,
  tabsize=2
}

\apptocmd{\frame}{}{\justifying}{}

\setbeamertemplate{navigation symbols}{}

\title{06 - JavaScript: sintaxe moderna (ES6+)}
\author{Prof. Alex Kutzke}
\date{Setembro de 2026}

\begin{document}

% ==========================================
\begin{frame}
	\begin{center}
		\large{Setor de Educação Profissional e Tecnológica}\\
		Tecnologia em Análise e Desenvolvimento de Sistemas\\
		\vspace{1cm}
		\small{DS122 - Desenvolvimento de Aplicações Web 1}
	\end{center}
	\titlepage
\end{frame}

% ==========================================
\begin{frame}{Objetivos da aula}
	Ao final desta aula você deve ser capaz de:
	\begin{enumerate}
		\item Executar uma expressão no console do navegador e ler a resposta;
		\item Ligar um arquivo \texttt{.js} externo com \texttt{defer} e confirmar
		      que ele executou;
		\item Escolher entre \texttt{const} e \texttt{let}, e dizer por que
		      \texttt{var} ficou de fora;
		\item Prever o resultado de uma comparação com \texttt{==} e com
		      \texttt{===};
		\item Escrever uma função nas três formas de hoje e converter uma na outra;
		\item Percorrer um array com \texttt{for...of} e com \texttt{forEach};
		\item Produzir arrays novos com \texttt{map} e \texttt{filter}, e um valor
		      único com \texttt{reduce};
		\item Representar o catálogo como array de objetos e extrair um relatório;
		\item Ler uma mensagem de erro do console e chegar à causa.
	\end{enumerate}
\end{frame}

% ==========================================
\begin{frame}{Roteiro}
	\tableofcontents
\end{frame}

% ==========================================
\section{Console}
% ==========================================

\begin{frame}[fragile]{O console do navegador}
	\texttt{F12}, aba \textbf{Console}. O mesmo painel do inspetor de CSS.

	\begin{minted}{javascript}
2 + 3
"Loja" + " " + "Exemplo"
document.title
document.querySelectorAll("article.produto").length
	\end{minted}

	\vspace{0.3cm}
	Uma linha, um resultado, sem recarregar nada.
\end{frame}

\begin{frame}{O que apareceu na última linha}
	\begin{itemize}
		\item o código tinha acesso ao conteúdo da página, sem nada salvo em
		      arquivo;
		\item \texttt{document} existia sem ter sido declarado;
		\item \texttt{querySelectorAll} e \texttt{length} vieram junto com ele.
	\end{itemize}

	\vspace{0.5cm}
	Esse acesso à página chama-se DOM, e é o assunto do próximo encontro.

	\vspace{0.3cm}
	Hoje o console é a calculadora onde o código é conferido antes de ir para o
	arquivo.
\end{frame}

\begin{frame}{Onde o JavaScript roda}
	Linguagem interpretada, de tipagem dinâmica, criada em 1995 para o navegador.

	\vspace{0.3cm}
	Cada navegador traz um motor: V8 (Chrome, Edge), SpiderMonkey (Firefox),
	JavaScriptCore (Safari).

	\vspace{0.3cm}
	Desde 2009 o V8 também roda fora do navegador, no Node.js. Neste semestre o
	servidor é PHP, e o JavaScript fica no cliente.

	\vspace{0.3cm}
	A sintaxe é a mesma nos dois lugares.
\end{frame}

\begin{frame}{ECMAScript e o que significa ES6}
	A especificação se chama \textbf{ECMAScript}. JavaScript é a implementação
	que os navegadores entregam.

	\vspace{0.4cm}
	\textbf{ES6}, ou ES2015, reorganizou a linguagem: \texttt{let}, \texttt{const},
	arrow functions, template strings, classes e promessas.

	\vspace{0.4cm}
	Desde 2015 sai uma versão por ano, com acréscimos menores.

	\vspace{0.4cm}
	Tutorial anterior a 2015 usa \texttt{var} e outras construções já
	substituídas. Reconhecer o estilo antigo evita copiá-lo.
\end{frame}

\begin{frame}[fragile]{O script na página}
	\begin{minted}{html}
<!-- 1. No atributo do elemento. Não use. -->
<button onclick="alert('oi')">Clique</button>

<!-- 2. Dentro da página. Só para teste rápido. -->
<script>
  console.log("rodou");
</script>

<!-- 3. Arquivo externo. É o que usamos. -->
<script src="js/app.js" defer></script>
	\end{minted}

	\vspace{0.2cm}
	Estrutura no HTML, aparência no CSS, comportamento no JavaScript.
\end{frame}

\begin{frame}[fragile]{Por que \texttt{defer}}
	Sem atributo nenhum, o navegador para de montar a página, baixa o arquivo e
	executa. O script roda antes de a página existir.

	\vspace{0.3cm}
	\texttt{defer}: baixa em paralelo, executa depois que a página está montada.

	\begin{minted}{html}
<head>
  <link rel="stylesheet" href="css/estilo.css">
  <script src="js/produtos.js" defer></script>
  <script src="js/app.js" defer></script>
</head>
	\end{minted}

	\vspace{0.2cm}
	Vários scripts com \texttt{defer} executam na ordem em que aparecem.
\end{frame}

\begin{frame}{Confirmar que o arquivo carregou}
	Caminho errado em \texttt{src} não aparece na tela: a página abre sem
	comportamento nenhum.

	\vspace{0.4cm}
	\begin{enumerate}
		\item a primeira linha do arquivo é um \texttt{console.log}, e a mensagem
		      aparece ao recarregar;
		\item na aba \textbf{Network}, existe uma linha para o arquivo com status
		      \texttt{200}.
	\end{enumerate}

	\vspace{0.4cm}
	Status \texttt{404} ali significa caminho errado, e ele é relativo ao arquivo
	HTML.
\end{frame}

% ==========================================
\section{Valores}
% ==========================================

\begin{frame}[fragile]{Tipagem dinâmica}
	O tipo pertence ao valor, não à variável.

	\begin{minted}{javascript}
typeof "Café"      // "string"
typeof 3.14        // "number"
typeof false       // "boolean"
typeof undefined   // "undefined"
typeof [1, 2, 3]   // "object"
typeof null        // "object"
	\end{minted}

	\vspace{0.2cm}
	Array é um tipo de objeto. Para perguntar se um valor é array:
	\texttt{Array.isArray(valor)}.

	\vspace{0.2cm}
	\texttt{typeof null} é defeito de 1995, nunca corrigido para não quebrar
	código existente.
\end{frame}

\begin{frame}[fragile]{\texttt{undefined} e \texttt{null}}
	\texttt{undefined} é o que o interpretador coloca. \texttt{null} é o que você
	coloca.

	\begin{minted}{javascript}
let preco;
console.log(preco);            // undefined

const produto = { nome: "Café" };
console.log(produto.desconto); // undefined

let selecionado = null;        // ninguém selecionado
	\end{minted}

	\vspace{0.2cm}
	Variável sem valor, propriedade que não existe e função sem \texttt{return}
	devolvem \texttt{undefined}.
\end{frame}

\begin{frame}[fragile]{\texttt{const} e \texttt{let}}
	\texttt{const} por padrão. \texttt{let} quando o valor for reatribuído.

	\begin{minted}{javascript}
const taxa = 0.1;
let total = 0;

total = total + 100;   // válido
taxa = 0.2;            // TypeError
	\end{minted}

	\vspace{0.3cm}
	Escrever \texttt{preco = 10} sem declarar cria uma global silenciosa.
\end{frame}

\begin{frame}[fragile]{\texttt{const} com array e objeto}
	\texttt{const} protege a ligação entre o nome e o valor, não o conteúdo.

	\begin{minted}{javascript}
const carrinho = ["café", "chá"];

carrinho.push("mate");   // válido
carrinho = ["outro"];    // TypeError
	\end{minted}

	\vspace{0.3cm}
	Array declarado com \texttt{const} continua aceitando \texttt{push},
	\texttt{pop} e atribuição por índice.
\end{frame}

\begin{frame}[fragile]{Escopo de bloco}
	Bloco é qualquer trecho entre chaves.

	\begin{minted}{javascript}
let x = "fora";

if (true) {
  let x = "dentro";
  console.log(x);     // dentro
}

console.log(x);       // fora
	\end{minted}
\end{frame}

\begin{frame}[fragile]{Por que \texttt{var} ficou de fora}
	\begin{minted}{javascript}
function testa() {
  for (var i = 0; i < 3; i++) {
    // ...
  }
  console.log(i);   // 3, o i sobreviveu ao for
}
	\end{minted}

	\vspace{0.2cm}
	\begin{itemize}
		\item \texttt{var} tem escopo de função, e vaza para fora do bloco;
		\item sofre \emph{hoisting}: lido antes da declaração, devolve
		      \texttt{undefined} em vez de erro.
	\end{itemize}

	\vspace{0.2cm}
	Aparece aqui para ser reconhecido em código alheio.
\end{frame}

\begin{frame}[fragile]{Template literals}
	Crase, com \texttt{\$\{\}} em volta de cada expressão.

	\begin{minted}{javascript}
const nome = "Café em grão";
const preco = 32.5;

// concatenação, o jeito antigo
console.log("O " + nome + " custa R$ " + preco + ".");

// template literal
console.log(`O ${nome} custa R$ ${preco}.`);

// qualquer expressão serve
console.log(`Com 10% off: R$ ${preco * 0.9}.`);
	\end{minted}
\end{frame}

\begin{frame}[fragile]{Strings}
	\begin{minted}{javascript}
const termo = "  Café Em Grão  ";

termo.trim()                       // "Café Em Grão"
termo.trim().toLowerCase()         // "café em grão"
"café em grão".includes("grão")    // true
"café em grão".startsWith("café")  // true
	\end{minted}

	\vspace{0.2cm}
	Strings são imutáveis: todo método devolve uma string nova, e o resultado
	precisa ser guardado ou usado ali mesmo.

	\vspace{0.2cm}
	Busca que ignora maiúsculas e minúsculas:

	\begin{minted}{javascript}
nome.toLowerCase().includes(busca.toLowerCase())
	\end{minted}
\end{frame}

\begin{frame}[fragile]{Números são ponto flutuante}
	\begin{minted}{javascript}
0.1 + 0.2        // 0.30000000000000004
	\end{minted}

	\vspace{0.2cm}
	O mesmo acontece em C, Java e Python: 0,1 não tem representação exata em
	base 2.

	\begin{minted}{javascript}
const total = 32.5 * 3;       // 97.5

total.toFixed(2)              // "97.50", string
Number(total.toFixed(2))      // 97.5

(1890).toLocaleString("pt-BR",
  { style: "currency", currency: "BRL" });
// "R$ 1.890,00"
	\end{minted}
\end{frame}

\begin{frame}[fragile]{Texto que vira número}
	Campo de formulário devolve sempre string.

	\begin{minted}{javascript}
"10" + 5        // "105", concatenação
"10" - 5        // 5, o - converte

Number("10")          // 10
Number("10,5")        // NaN
Number("")            // 0
parseFloat("10.5 reais")  // 10.5
parseInt("42px")          // 42
	\end{minted}

	\vspace{0.2cm}
	\texttt{NaN} é do tipo \texttt{number} e não é igual a si mesmo.

	\vspace{0.1cm}
	Teste com \texttt{Number.isNaN(valor)}.
\end{frame}

% ==========================================
\section{Fluxo}
% ==========================================

\begin{frame}[fragile]{\texttt{===} e \texttt{==}}
	\texttt{===} compara valor e tipo. \texttt{==} converte antes de comparar.

	\begin{minted}{javascript}
1 == "1"        // true
1 === "1"       // false

0 == false      // true
0 === false     // false

null == undefined    // true
null === undefined   // false
	\end{minted}

	\vspace{0.2cm}
	Use \texttt{===} e \texttt{!==} sempre.
\end{frame}

\begin{frame}[fragile]{Os seis valores falsos}
	\begin{minted}{javascript}
false
0
""          // string vazia
null
undefined
NaN
	\end{minted}

	\vspace{0.2cm}
	Todo o resto vira \texttt{true}, inclusive \texttt{"0"}, \texttt{"false"},
	\texttt{[]} e \texttt{\{\}}.

	\vspace{0.2cm}
	Array vazio dentro de um \texttt{if} entra no bloco.
\end{frame}

\begin{frame}[fragile]{Decisão}
	\begin{minted}{javascript}
if (preco > 100) {
  console.log("caro");
} else if (preco > 20) {
  console.log("médio");
} else {
  console.log("barato");
}

const frete = preco > 100 ? 0 : 15;

if (preco >= 20 && preco <= 50) { /* ... */ }
	\end{minted}

	\vspace{0.2cm}
	Operadores lógicos: \texttt{\&\&}, \texttt{||} e \texttt{!}.
\end{frame}

\begin{frame}[fragile]{Repetição}
	\begin{minted}{javascript}
const frutas = ["café", "chá", "mate"];

// quando o índice importa
for (let i = 0; i < frutas.length; i++) {
  console.log(i, frutas[i]);
}

// quando só o valor importa
for (const fruta of frutas) {
  console.log(fruta);
}
	\end{minted}

	\vspace{0.2cm}
	O contador é \texttt{let} porque é reatribuído. A variável do
	\texttt{for...of} recebe valor novo a cada volta, e pode ser \texttt{const}.
\end{frame}

% ==========================================
\section{Funções}
% ==========================================

\begin{frame}[fragile]{Três formas de escrever}
	\begin{minted}{javascript}
// declaração
function calculaTotal(preco, quantidade) {
  return preco * quantidade;
}

// expressão de função
const calculaTotal = function (preco, quantidade) {
  return preco * quantidade;
};

// arrow function (ES6)
const calculaTotal = (preco, quantidade) => {
  return preco * quantidade;
};
	\end{minted}
\end{frame}

\begin{frame}[fragile]{Arrow function, formas curtas}
	\begin{minted}{javascript}
// corpo de uma expressão: chaves e return somem juntos
const calculaTotal = (preco, qtd) => preco * qtd;

// um parâmetro: parênteses opcionais
const dobro = valor => valor * 2;

// nenhum parâmetro: parênteses obrigatórios
const agora = () => new Date();
	\end{minted}
\end{frame}

\begin{frame}[fragile]{A armadilha do retorno implícito}
	\begin{minted}{javascript}
const dobro = valor => { valor * 2 };

dobro(4);    // undefined
	\end{minted}

	\vspace{0.3cm}
	Com chaves, o \texttt{return} precisa estar escrito. Sem chaves, o valor da
	expressão é devolvido.

	\vspace{0.3cm}
	É o defeito mais comum de quem está começando, e ele não gera erro nenhum no
	console.
\end{frame}

\begin{frame}[fragile]{Parâmetro com valor padrão}
	\begin{minted}{javascript}
const calculaTotal = (preco, quantidade = 1) =>
  preco * quantidade;

calculaTotal(32.5);      // 32.5
calculaTotal(32.5, 3);   // 97.5
	\end{minted}

	\vspace{0.2cm}
	Chamar com menos argumentos do que os parâmetros não é erro: os que faltam
	ficam \texttt{undefined}, e a conta devolveria \texttt{NaN}.
\end{frame}

\begin{frame}[fragile]{Função é um valor}
	\begin{minted}{javascript}
const aplicar = (valor, operacao) => operacao(valor);

aplicar(4, dobro);               // 8
aplicar(4, valor => valor + 1);  // 5
	\end{minted}

	\vspace{0.3cm}
	Função guardada em variável, passada como argumento, devolvida por outra
	função.

	\vspace{0.3cm}
	Função passada como argumento chama-se \emph{callback}, e é o que torna
	possível o \texttt{forEach} da próxima seção.
\end{frame}

% ==========================================
\section{Arrays e objetos}
% ==========================================

\begin{frame}[fragile]{Array}
	\begin{minted}{javascript}
const precos = [32.5, 18, 74.9, 12];

precos[0]                   // 32.5
precos.length               // 4
precos[precos.length - 1]   // 12
precos[10]                  // undefined

precos.push(50);      // acrescenta no fim
precos.pop();         // remove do fim
precos.unshift(5);    // acrescenta no começo
precos.shift();       // remove do começo
	\end{minted}
\end{frame}

\begin{frame}[fragile]{Percorrer: \texttt{forEach}}
	\begin{minted}{javascript}
const produtos = ["Café", "Chá", "Mate"];

produtos.forEach(produto => console.log(produto));

produtos.forEach((produto, indice) => {
  console.log(`${indice}: ${produto}`);
});
	\end{minted}

	\vspace{0.2cm}
	Não devolve nada e não pode ser interrompido: \texttt{break} não funciona
	dentro dele. Para parar antes do fim, \texttt{for...of}.
\end{frame}

\begin{frame}[fragile]{Transformar: \texttt{map}}
	Devolve um array \textbf{novo}, do mesmo tamanho.

	\begin{minted}{javascript}
const precos = [32.5, 18, 74.9];

const comDesconto = precos.map(preco => preco * 0.9);
// [29.25, 16.2, 67.41000000000001]

precos;   // [32.5, 18, 74.9], inalterado
	\end{minted}

	\vspace{0.2cm}
	A cauda de casas do terceiro valor é o ponto flutuante da seção anterior. Ela
	some na exibição, com \texttt{toFixed} ou \texttt{toLocaleString}.
\end{frame}

\begin{frame}[fragile]{Selecionar: \texttt{filter}}
	Devolve um array novo com os elementos que passaram no teste.

	\begin{minted}{javascript}
const precos = [32.5, 18, 74.9, 12];

const baratos = precos.filter(preco => preco < 30);
// [18, 12]
	\end{minted}

	\vspace{0.3cm}
	A função precisa devolver \texttt{true} ou \texttt{false}. Escrever a condição
	entre chaves e esquecer o \texttt{return} devolve array vazio, porque
	\texttt{undefined} é valor falso.
\end{frame}

\begin{frame}[fragile]{Resumir: \texttt{reduce}}
	\begin{minted}{javascript}
const precos = [32.5, 18, 74.9, 12];

const soma = precos.reduce(
  (acumulado, preco) => acumulado + preco,
  0
);
// 137.4
	\end{minted}

	\vspace{0.2cm}
	\texttt{acumulado} é o que voltou da chamada anterior; na primeira volta, é o
	segundo argumento.

	\vspace{0.2cm}
	Escreva sempre o valor inicial.
\end{frame}

\begin{frame}[fragile]{Procurar}
	\begin{minted}{javascript}
const precos = [32.5, 18, 74.9];

precos.includes(18);                    // true
precos.find(preco => preco > 30);       // 32.5
precos.findIndex(preco => preco > 30);  // 0
precos.some(preco => preco > 70);       // true
precos.every(preco => preco > 10);      // true
	\end{minted}

	\vspace{0.2cm}
	\texttt{find} devolve o elemento, ou \texttt{undefined}. \texttt{filter}
	devolve sempre um array.
\end{frame}

\begin{frame}[fragile]{Encadear}
	\begin{minted}{javascript}
const total = precos
  .filter(preco => preco < 50)   // [32.5, 18, 12]
  .map(preco => preco * 0.9)     // [29.25, 16.2, 10.8]
  .reduce((soma, p) => soma + p, 0);
// 56.25
	\end{minted}

	\vspace{0.2cm}
	\texttt{map} e \texttt{filter} devolvem arrays, então as chamadas se
	encadeiam e a leitura é da esquerda para a direita.
\end{frame}

\begin{frame}[fragile]{Ordenar: \texttt{sort}}
	\texttt{sort} altera o array original e, sem argumento, ordena como texto.

	\begin{minted}{javascript}
const precos = [32.5, 18, 74.9, 120];

precos.sort();               // [120, 18, 32.5, 74.9]
precos.sort((a, b) => a - b); // [18, 32.5, 74.9, 120]
precos.sort((a, b) => b - a); // decrescente
	\end{minted}

	\vspace{0.2cm}
	A comparação devolve negativo quando \texttt{a} vem antes, positivo quando
	\texttt{b} vem antes, zero quando tanto faz.
\end{frame}

\begin{frame}[fragile]{Objeto}
	\begin{minted}{javascript}
const produto = {
  id: 1,
  nome: "Café em grão",
  categoria: "Bebidas",
  preco: 32.5,
  desconto: 0.15
};

produto.nome        // "Café em grão"
produto["nome"]     // igual

const campo = "preco";
produto[campo]      // 32.5
	\end{minted}

	\vspace{0.2cm}
	Colchete quando o nome da chave vem de uma variável.
\end{frame}

\begin{frame}[fragile]{Array de objetos}
	É o formato em que os dados chegam de uma API, e o do seu catálogo.

	\begin{minted}{javascript}
const produtos = [
  { id: 1, nome: "Café em grão",  preco: 32.5 },
  { id: 2, nome: "Chá de hibisco", preco: 18  },
  { id: 3, nome: "Mel silvestre",  preco: 24.9 }
];

produtos.map(p => p.nome);
produtos.filter(p => p.preco < 30);
produtos.find(p => p.id === 3);
produtos.reduce((s, p) => s + p.preco, 0);
	\end{minted}
\end{frame}

\begin{frame}[fragile]{Desestruturação e \texttt{console.table}}
	\begin{minted}{javascript}
const { nome, preco } = produto;

const linha = ({ nome, preco }) =>
  `${nome}: R$ ${preco.toFixed(2)}`;

console.table(produtos);
	\end{minted}

	\vspace{0.3cm}
	\texttt{console.table} monta uma linha por item e uma coluna por campo. É a
	forma mais rápida de conferir se o filtro selecionou o que deveria.
\end{frame}

% ==========================================
\section{Erros}
% ==========================================

\begin{frame}[fragile]{A mensagem no console}
	Diferente do CSS, que ignora em silêncio, o JavaScript interrompe e escreve.

	\begin{minted}{text}
Uncaught TypeError: produtos.filtrar is not a function
    at app.js:14:20
	\end{minted}

	\vspace{0.3cm}
	\begin{tabular}{ll}
		\toprule
		\texttt{SyntaxError}    & chave ou aspas sem fechar; nada roda \\
		\texttt{ReferenceError} & nome não declarado, ou digitado errado \\
		\texttt{TypeError}      & o valor existe, mas não é do tipo pedido \\
		\bottomrule
	\end{tabular}
\end{frame}

\begin{frame}[fragile]{O erro mais comum}
	\begin{minted}{javascript}
const produto = produtos.find(p => p.id === 99);
console.log(produto.nome);
// TypeError: Cannot read properties of undefined
//   (reading 'nome')
	\end{minted}

	\vspace{0.3cm}
	O \texttt{find} não encontrou e devolveu \texttt{undefined}. O erro só
	apareceu na linha seguinte, quando alguém pediu \texttt{.nome} daquilo.

	\vspace{0.3cm}
	A causa costuma estar uma linha acima da linha apontada.
\end{frame}

\begin{frame}{Quando não há erro nenhum}
	Resultado errado, console limpo. Três conferências:

	\begin{enumerate}
		\item \texttt{console.log} do valor antes e depois da operação suspeita;
		\item \texttt{typeof valor}, para descobrir que o número era uma string;
		\item painel \textbf{Sources}, ponto de parada na linha, página
		      recarregada: a execução para ali e o painel mostra cada variável.
	\end{enumerate}

	\vspace{0.4cm}
	O terceiro item é o depurador, e ele volta com mais uso quando entrarmos em
	eventos.
\end{frame}

\begin{frame}{Parte prática}
	Sobre as suas páginas, ou sobre o pacote de exemplo da aula:

	\begin{enumerate}
		\item criar \texttt{js/app.js} e ligá-lo com \texttt{defer};
		\item escrever \texttt{js/produtos.js} com o catálogo em array de objetos;
		\item formatar preço e imprimir uma linha por produto;
		\item aplicar desconto;
		\item busca por nome e filtro por faixa de preço, combináveis;
		\item relatório com \texttt{reduce};
		\item quebrar de propósito e ler o erro.
	\end{enumerate}

	\vspace{0.3cm}
	Nada disso muda a página. Levar o resultado para a tela é o próximo encontro.
\end{frame}

\begin{frame}{Tarefa}
	Os exercícios valem nota e compõem o item \emph{Exercícios em sala}.

	\vspace{0.3cm}
	Enunciado no \texttt{README.md} do repositório-modelo; endereço e prazo na
	UFPR Virtual.

	\vspace{0.3cm}
	Entrega pelo \emph{fork} no GitLab, individual ou em dupla, sem uso de IA
	generativa para produzir o código.

	\vspace{0.5cm}
	A busca e o filtro de hoje são dois requisitos da Entrega 2 do trabalho
	prático, que vence na aula anterior à Prova 2.
\end{frame}

\end{document}
